\documentclass[a4paper,12pt, noSujetB]{article}

\usepackage{layout_evaluation}

\classe{Seconde}
\titre{Triangles de Pick - arithmétique}
\evaltype{Travail en groupe}
\evalnumber{1}
\date{Toussaint}

\begin{document}
    \section*{Consignes:}

    Avec votre groupe, choisissez un des exercices. Ce sont des exercices volontairement difficiles et il n'est pas attendu que vous répondiez à toutes les questions. Il est possible d'avoir une très bonne note en traitant seulement 3 questions sur 10.

    \medskip
    A la rentrée, certains groupes présenteront à l'oral leur travail. C'est votre capacité à présenter un exercice devant la classe qui est noté et non pas la quantité de questions auxquelles vous avez répondu. 
    
    \medskip
    Il est demandé de préparer un support visuel pour vous aider dans votre présentation. Vous pouvez par exemple scanner votre copie, ou préparer des slides. Vous pouvez utiliser le tableau lors de votre présentation. 

    \medskip
    Tous les groupes rendront une copie, qui sera corrigé et non noté.
   
    \medskip
    Utiliser une intelligence artificielle pour se travail vous desservira pour les raisons suivantes: 
        \begin{itemize}
            \item Ce n'est pas la quantité de questions traitée qui est prise en compte mais votre présentation orale.
            \item Cet exercice doit faire travailler vos neurones, pas celles d'un serveur en Californie.
            \item La notation (cf. ci-dessous) est adaptée pour prendre en compte la démarche personnelle dans la résolution des problèmes. Il est probable qu'un modèle propose une approche générique et que vous ne la compreniez pas suffisament pour répondre aux questions.
        \end{itemize}


        \section*{Grille de notation}
        \begin{itemize}
            \item La durée est respectée ($\sim$ 10 min).
            \item Tous les membres du groupe ont eu la paroles.
            \item Les explications orales sont claires.
            \item Le vocabulaire mathématique est correctement prononcé.
            \item Le contenu est mathématiquement correct.
            \item La démarche de résolution est originale et bien comprise.
            \item Le support visuel est lisible et soigné.
            \item Le support visuel aide à la compréhension, le tableau est utilisé (si besoin).
            \item Le groupe sait répondre aux questions et justifier sa démarche.
        \end{itemize}

\end{document}